\section{表格}
\subsection{三线表与表头}
\begin{frame}[fragile]{三线表与粗体表头}
  \begin{tabular}{@{}lrr@{}}
    \toprule
    \tablehead{方法} & \tablehead{准确率} & \tablehead{耗时（秒）} \\
    \midrule
    基线 & 0.81 & 12.3 \\
    改进 & 0.87 & 10.8 \\
    \bottomrule
  \end{tabular}
  \begin{verbatim}
\tablehead{方法} & \tablehead{准确率} \\
  \end{verbatim}
  \slideNote{表头每个单元格用 tablehead 标记，统一使用衬线粗体。
  l、c、r 分别控制左、中、右对齐；三线表使用 top/mid/bottomrule。}
\end{frame}

\subsection{自适应宽度与合并表头}
\begin{frame}{自动列宽与多行内容}
  \small
  \begin{tabularx}{\linewidth}{@{}lYY@{}}
    \toprule
    \tablehead{方法} & \tablehead{适用场景} & \tablehead{注意事项} \\
    \midrule
    文献调研 & 归纳不同工作的共同问题、方法与发展趋势 & 统一比较口径，注明资料的时间范围 \\
    系统调研 & 检查公开实现与实验设置，进行复现 & 记录版本、运行环境和关键参数 \\
    \bottomrule
  \end{tabularx}
  \sourceNote{示例内容；Y 列自动分配宽度并换行，文字左对齐。}
  \slideNote{表格行距统一为 1.2。使用 tabularx 与 Y 列处理长文字，
  不需要手工在每个单元格中插入换行。}
\end{frame}

\begin{frame}{合并表头与表格引用}
  \begin{table}
    \centering
    \small
    \begin{tabular}{@{}lrrrr@{}}
      \toprule
      & \multicolumn{2}{c}{\tablehead{数据集 A}} & \multicolumn{2}{c}{\tablehead{数据集 B}} \\
      \cmidrule(lr){2-3}\cmidrule(l){4-5}
      \tablehead{方法} & \tablehead{均值} & \tablehead{标准差} & \tablehead{均值} & \tablehead{标准差} \\
      \midrule
      基线 & 0.81 & 0.03 & 0.79 & 0.04 \\
      改进 & 0.87 & 0.02 & 0.85 & 0.03 \\
      \bottomrule
    \end{tabular}
    \caption{两组数据集上的示例结果}\label{tab:demo}
  \end{table}
  表~\ref{tab:demo} 演示 multicolumn 合并单元格与 cmidrule 分组横线。
  \sourceNote{数值为排版示例，不代表实际实验结果。}
\end{frame}
\begin{frame}{分栏与表格}
  \begin{columns}[T,onlytextwidth]
    \column{0.50\textwidth}
      \textbf{要点}（左栏）
      \begin{enumerate}
        \item 栏宽按内容分配，两栏合计不超过 $1$。
        \item \texttt{[T]} 让两栏顶部对齐；\texttt{onlytextwidth} 避免整体溢出。
        \item 栏内可放列表、公式或图片。
      \end{enumerate}
    \column{0.46\textwidth}
      \small
      \begin{tabularx}{\linewidth}{@{}lY@{}}
        \toprule
        \tablehead{维度} & \tablehead{说明} \\
        \midrule
        数据 & 填写数据来源与规模 \\
        指标 & 填写评价指标与口径 \\
        结论 & 填写该维度的结论 \\
        \bottomrule
      \end{tabularx}
      \sourceNote{表格使用 \texttt{tabularx} 与 \texttt{Y} 列自适应宽度。}
  \end{columns}
\end{frame}
